\section{The exact remaining geometric criterion and the actual determinant}
\label{sec:exact-geometric-gate}

The preceding partial-strip theorem allows the remaining question to
be stated for the literal cutoff operators at every polynomial weight.
It no longer has an unproved trace-class or boundary-closure premise.
This section makes that consequence explicit and records the actual
Fredholm determinant of the same quotient.

Use the rational trivial sector, put $a=1/2$, and fix
$\alpha>0$, $s\ge0$. Let $r_\rho=r_\rho(\alpha,s)$ be
(MA-WP.3) for $u=-\alpha$, $v=\alpha$.
For $g\in C_c^\infty(\mathbb R)$ with $g(0)=0$, write
$G(z)=\mathcal Mg(z)$ and
\[
\begin{aligned}
 A_0(g)&=\int_0^\infty
   \frac{e^{-x/2}}{1-e^{-2x}}\bigl(g(x)+g(-x)\bigr)\,dx,\\
 B(g)&=\sum_{p,k\ge1}(\log p)p^{-k/2}
                 \bigl(g(k\log p)+g(-k\log p)\bigr),\\
 P(g)&=G(1/2)+G(-1/2),\qquad
 \mathcal W(g)=A_0(g)+B(g)-P(g),\\
 \mathcal G_{\alpha,s}(g)&=
      \lim_R\operatorname{Tr}(M_{\chi_R}C_gP_G).
\end{aligned}
\tag{MA-EG.1}
\]
The cutoffs satisfy the exact hypotheses of MA-WP-003.
The limit exists unconditionally by that theorem, and does not
depend on their particular shape or exhaustion rate.
The prime sum in (MA-EG.1) is finite. The archimedean integral
is ordinary convergent at zero because $g(0)=0$.

\paragraph{MA-EG-001 (criterion for the actual unchanged geometric identity).}
For these actual operators,
\[
 \mathcal G_{\alpha,s}(g)-\mathcal W(g)
    =\sum_\rho (m_\rho-r_\rho)G(\rho-1/2).
 \tag{MA-EG.2}
\]
In particular, the following are equivalent:
\begin{enumerate}
\item $\mathcal G_{\alpha,s}(g)=\mathcal W(g)$ for every compact
smooth $g$ with $g(0)=0$;
\item $r_\rho=m_\rho$ for every nontrivial zero;
\item no zero lies strictly outside the closed centered strip,
and every zero on its boundary has $m_\rho\le b_s$,
where $b_s$ is (MA-WP.2).
\end{enumerate}
If a zero fails condition~2, a compact smooth witness with
$g(0)=0$ and nonzero difference in (MA-EG.2) exists.
Requiring condition~1 for every $\alpha>0$, with any fixed
$s\ge0$, is exactly the Riemann zero-location statement.
This is an equivalence with an unproved independent-method step,
not a new zero-location proof.

\begin{proof}
The classical full explicit formula in the normalization of
MA-TR-2--5, equivalently Meyer's Theorem~5.8 and its spectral
identification Theorem~5.11, gives
\[
             \mathcal W(g)=-\sum_\rho m_\rho G(\rho-1/2).
\]
MA-WP-003 gives the actual operator limit
$\mathcal G_{\alpha,s}(g)=-\sum_\rho r_\rho G(\rho-1/2)$.
Both sums converge absolutely. Subtraction proves (MA-EG.2).
The equivalence of conditions~2 and~3 is exactly (MA-WP.3),
and condition~2 implies condition~1.

For the converse, the two spectral functionals extend continuously
to $\mathcal E$ and annihilate the centered smooth ideal
$I_{1/2}$: the local zero count and rapid vertical Mellin decay
prove continuity, and every ideal element has all the required
zero jets. The restricted-trace criterion MT-M30, (M30.13)--(M30.14),
therefore applies with weights $r_\rho$.
Here is its explicit fixed-sector witness, to make clear what
happens at a lost boundary multiplicity. If
$m_{\rho_0}-r_{\rho_0}>0$, MT-M18 gives $f\in\mathcal E$
whose centered Mellin value is one at $\rho_0$ and zero at every
other zero, with the prescribed higher jets. Choose a nonzero
$\theta\in I_{1/2}$ and put
\[
 b=\frac{\theta*\theta^*}{\|\theta\|_2^2},\qquad
 \theta^*(x)=\overline{\theta(-x)}.
\]
The ideal property gives $b\in I_{1/2}$ and direct evaluation
gives $b(0)=1$. Thus $f_0=f-f(0)b$ vanishes at zero and has
exactly the same spectral values as $f$. Choose a fixed
$\kappa\in C_c^\infty$ equal to one near zero and put
$\kappa_R(x)=\kappa(x/R)$. Then
$\kappa_R f_0\to f_0$ in $\mathcal E$. Its spectral difference
tends to the nonzero number $m_{\rho_0}-r_{\rho_0}$.
For all sufficiently large $R$, $g=\kappa_R f_0$ is the required
compact witness. This proves the converse and the witness assertion.
The construction is conditional on an actually omitted zero;
it does not assert that such a zero exists for $\zeta$.

If all zeros are on the central line, every positive-width
strip contains them strictly, regardless of their multiplicities.
Conversely, any zero with $|\operatorname{Re}\rho-1/2|>0$
is exterior for a smaller positive width, contradicting
condition~2 for that width. This proves the final equivalence.
\end{proof}

The same proof works in each fixed Dirichlet or finite-order Hecke
character sector, with its own normalized full Weil distribution
from MA-CH-001 and its own pole term. The full-family norm-one
splitting is precisely MT-M30. Thus the actual cutoff identity in
all rational Dirichlet strips follows using the specified theorem
$A$. Using that theorem to remove the discarded-divisor term is
a zero-location input, hence circular for the requested independent
method test. For the non-induced imaginary-quadratic character of
MA-CH-003, the corresponding Hecke zero-location statement is still
the condition; $M$ supplies it, whereas $A$ does not.

\paragraph{MA-EG-002 (Fredholm determinant of the actual partial quotient).}
In the general fixed-character setting of MA-WP, for every
$g\in\mathcal E$ the actual trace-class quotient convolution has
\[
 \det\bigl(1-wC_g^Q\bigr)
       =\prod_{\Xi(\rho)=0}
           \bigl(1-w\mathcal Mg(\rho-a)\bigr)^{r_\rho},
 \qquad w\in\mathbb C.
 \tag{MA-EG.3}
\]
The product converges locally uniformly. Its derivative at zero
is the negative of the actual trace (MA-WP.20).
It retains every surviving finite-block length, including the exact
boundary truncations, without an assumption of simple zeros or of
a positive operator.

\begin{proof}
Let $A_{n,g}$ be the actual finite-rank operators (MA-WP.13).
Their finite sums $T_N=\sum_{|n|\le N}A_{n,g}$ converge to
$C_g^Q$ in trace norm by (MA-WP.15).
The mutually annihilating bounded projections $P_n$ give a finite
invariant direct sum $\bigoplus_{|n|\le N}V_n$, and $T_N$ is
zero on the kernel of the bounded projection
$\sum_{|n|\le N}P_n$. On a retained local factor
$\mathbb C[z]/(z^{r_\rho})$, multiplication by the Taylor class
of $\mathcal Mg$ is triangular with diagonal
$\mathcal Mg(\rho-a)$. Finite-dimensional determinant theory gives
\[
 \det(1-wT_N)=
   \prod_{\rho\text{ in the blocks }|n|\le N}
        \bigl(1-w\mathcal Mg(\rho-a)\bigr)^{r_\rho}.
\]
Trace-norm continuity of the Fredholm determinant now gives
convergence of the left side locally uniformly in $w$.
An exact reference, including the needed quantitative estimate,
is Kostenko, \emph{Trace Ideals with Applications},
Theorem~3.4.4 and (3.4.19), printed pp.~39--40:
\[
 |\det(1+A)-\det(1+B)|
       \le\|A-B\|_1\exp(1+\|A\|_1+\|B\|_1).
\]
The right-side products converge locally uniformly because
$\sum_\rho r_\rho|\mathcal Mg(\rho-a)|<\infty$, as proved
in MA-WP-002. Taking the limits proves (MA-EG.3).
Differentiation at zero is justified by local uniform convergence
and gives the displayed trace. No claim that the finite projections
are orthogonal has entered the argument.
\end{proof}

\paragraph{Grothendieck placement and the precise method boundary.}
The retained local algebra has length $r_\rho$ over
$\mathbb C$ and defines a finite-length module supported at the
spectral point $\rho-a$ on the complex affine line. Its class in
the Grothendieck group $K_0$ of the abelian category of finite-length
$\mathbb C[z]$-modules records this length at that point.
Exterior powers over $\mathbb C$ of the underlying finite-dimensional
vector spaces, with their induced endomorphisms, produce the finite
block determinants used above. Formula (MA-EG.3) is their nuclear
Fredholm limit in the actual completion. This is the cote~1
comparison at scans~079--080 and~111--112, via
ALPOGE\_BRIDGES BR-TR-01/02 and F-SM-11/12, with the actual
operator rather than only an associated graded used here.
The scan locators are inherited, not newly inspected in this run.
The determinant and trace do not recover the nilpotent extension
structure within $\mathbb C[z]/(z^{r_\rho})$.
No scheme over $\mathbb Z$, arithmetic cohomology, or
$\ell$-adic weight realization is constructed by this analytic
change of completion. In particular it supplies no faithful
full-jet positivity contrary to MT-M20/M23.

The source bridge remains the identification from the unfinished
\S12.3 on Masullo p.~514 to the global use of the local theorem
in (12.5.1.4), p.~520. MA-EG-001 identifies exactly what the
corrected actual objects require at that bridge. It does not
turn a missing proof into a mathematical counterexample under
$A$, and it does not infer nonexistence of every possible new
unbounded intertwiner. The direct bounded inverse obstruction,
the proved conditional full bridge, and the concrete regularized
commutator have their separate scopes in MA-MO, MA-HM and MA-RS.

\paragraph{Source details.}
Kostenko, A. (n.d.). \emph{Trace ideals with applications}
[Lecture notes], Theorem~3.4.4 and (3.4.19), pp.~39--40.
\url{https://users.fmf.uni-lj.si/kostenko/teach/IdealsNotes.pdf}.
The cited theorem and its complete proof were read in the
author-hosted version on 21 November 2026; no unverified publication
year is assigned. Meyer (2005), Theorem~5.8, author pp.~43--44,
and Theorem~5.11, p.~47, are used with the exact normalizations
already checked in MA-MEY. The inherited restricted-test theorem
is GRH\_NIGHT MT-M30, (M30.13)--(M30.14); the new input is
MA-WP's unconditional identification of the literal weighted
quotient and cutoff trace.
